\shorthandoff{<>}% Babel español trata como abreviaturas los delimitadores de curvatura de Xy-pic.

\begin{corollaire}
\label{XIII.3.2}
Sean $k$ un cuerpo de característica $p\geq 0$, $p'$ el conjunto de los
números primos distintos de~$p$ y $f\colon X\to k$ un morfismo coherente.
\begin{enumerate}
\item[1)] Para todo haz de conjuntos~$F$, $(F,f)$ es cohomológicamente
propio en dimensión $\leqslant 0$.
\item[2)] Supongamos que se cumple una de las dos condiciones siguientes:
\begin{enumerate}
\item[a)] $f$ es de tipo finito y los esquemas de tipo finito de dimensión
$\leqslant\dim X$ sobre una clausura algebraica de~$k$ son fuertemente
desingularizables.
\item[b)] Los esquemas de tipo finito sobre una clausura algebraica de~$k$
son
\marginpar{410}
fuertemente desingularizables.
\end{enumerate}
Entonces, para todo haz de ind-$p'$-grupos~$F$, $(F,f)$ es
cohomológicamente propio en dimensión $\leqslant 1$.
\end{enumerate}
\end{corollaire}

Sea $F$ un haz de conjuntos (\resp de ind-$p'$-grupos). Según
SGA~4 IX~2.7.2, puede escribirse~$F$ como límite inductivo filtrante
$$
F=\varinjlim F_i,
$$
donde los~$F_i$ son haces de conjuntos (\resp de ind-$p'$-grupos)
constructibles. Como $f$ es coherente, $f_*$ (\resp $\R^1f_*$) conmuta
con los límites inductivos (SGA~4 VII~3.3). Si se sabe que los~$(F_i,f)$
son cohomológicamente propios en dimensión $\leqslant 0$ (\resp que todo
haz de conjuntos es cohomológicamente propio en dimensión $\leqslant 0$
y que los~$(F_i,f)$ son cohomológicamente propios en dimensión
$\leqslant 1$), lo mismo valdrá para~$(F,f)$. Puede suponerse, por tanto,
que~$F$ es constructible.

Si se supone que~$f$ es de tipo finito (\resp que satisface la condición
a)), el resultado se sigue de~\Ref{XIII.3.1}~1)\kern2pt b) (\resp de
\Ref{XIII.3.1}~2)\kern2pt b)). Demostremos ahora~\Ref{XIII.3.2} sin
suponer ya que~$f$ sea de tipo finito. Para todo esquema~$S'$ sobre~$k$
y todo punto geométrico~$\overline{s}$ de~$S'$, se denota por
$\overline{k}$ (\resp $\overline{S'}$) la localización estricta de~$k$
en~$\overline{s}$ (\resp de~$S'$ en~$\overline{s}$), por~$\overline{X}$
la imagen inversa de~$X$ sobre~$\overline{k}$, y se considera el cuadrado
cartesiano
$$
\xymatrix@C=1.2cm{
\overline{X}\ar[d]_-{\overline{f}}&\ar[l]_g \overline{X'}\ar[d]^-{\overline{f}'}\\ \overline{k}&\ar[l]\overline{S'}
}
$$
Basta demostrar que, para todo~$S'$ y todo~$\overline{s}$, el morfismo
canónico
$$
\H^0(\overline{X},\overline{F})\to
\H^0(\overline{X'},\overline{F'})\qquad
\text{(\resp $\H^1(\overline{X},\overline{F})\to
\H^1(\overline{X'},\overline{F'})$)}
$$
es un isomorfismo. Basta demostrar que se tienen las relaciones
\marginpar{411}
\begin{equation*}
\label{eq:XIII.3.*}
\tag{$*$} \overline{F}\simeq g_*g^*\overline{F}\qquad
\ifthenelse{\boolean{orig}}
{\text{(\resp $\R^1g_*(g^*\overline{F}))=0$).}}
{\text{(\resp $\R^1g_*(g^*\overline{F})=0$).}}
\end{equation*}
Ahora bien, en esta forma la cuestión es local sobre~$X$ para la topología
étale. Puede suponerse, por tanto, que~$X$ es afín; por paso al límite,
puede suponerse que~$X$ es de tipo finito sobre~$k$. Se sabe entonces que
$(F,f)$ es cohomológicamente propio en dimensión $\leqslant 0$ (\resp
$\leqslant 1$), y que lo mismo ocurre al reemplazar~$X$ por un esquema
étale de tipo finito sobre~$X$, lo que demuestra~\eqref{eq:XIII.3.*}.
